\documentclass[11pt]{article}
\usepackage{amsmath, amssymb, geometry}
\usepackage[UTF8,fontset=fandol]{ctex} % compile with xelatex for CJK support and fonts
\geometry{margin=1in}
\title{Phase-Field Crystal Framework for Single-Crystal Copper: Summary}
\author{Project Memo}
\date{\today}
\begin{document}
\maketitle

\section{Background and Objectives}
The project replicates the ductile-fracture framework of Xu \emph{et al.}\cite{paper} to predict the tensile response of single-crystal copper (Cu) informed by EBSD observations. The primary objectives are: (i) build synthetic or EBSD-derived $[111]$ Cu structures with defects, (ii) couple Phase-Field Crystal (PFC) density evolution with phase-field fracture and crystal plasticity, (iii) solve mechanical equilibrium for anisotropic FCC copper, evolve plasticity and cracks, and export results for Ovito/ParaView, and (iv) compare simulated crack paths, stress--strain curves, and fatigue metrics (Paris, Coffin--Manson) with experiments.

\section{Energy Functional and Constraints}
State variables include the PFC density $\psi$, crack field $\phi\in[0,1]$, plastic strain $\varepsilon^p$, and grain/orientation markers $\eta_g$. The free energy reads
\begin{equation}
F = \int_{\Omega} \left[(1-\phi)^2 f_{\text{elas}}^{+} + f_{\text{elas}}^{-} + f_{\text{frac}} + f_{\text{grain}} + f_{\text{PFC}} \right] \, d\Omega.
\end{equation}
Elastic energy employs the cubic stiffness tensor $C_{ijkl}$, rotated to the local orientation $R$ via $C'_{ijkl}=R_{ip}R_{jq}R_{kr}R_{ls} C_{pqrs}$. Spectral decomposition of the strain tensor provides tensile/compressive splits: $\varepsilon = Q\Lambda Q^{\top}$, and $\Lambda^{+}$ keeps positive eigenvalues.
Here $\psi$ is the dimensionless PFC density deviation, $\phi$ is the damage/crack variable, $\varepsilon^p$ is the plastic strain, $\eta_g$ labels grain orientation, $f_{\text{elas}}^{\pm}$ are the positive/negative parts of the elastic energy density built from $\varepsilon - \varepsilon^p$, and $f_{\text{PFC}}$ captures lattice-scale ordering.

Fracture energy
\begin{equation}
f_{\text{frac}} = G_c(\eta_g, \varepsilon_{\text{eq}}) \left[ \frac{\phi^2}{2\ell_0} + \frac{\ell_0}{2} |\nabla \phi|^2 \right]
\end{equation}
uses a toughness degradation law
\begin{equation}
G_c = G_{c0} \left(0.5 - g_{\text{res}} \tanh \left[ 2 \left(\frac{\varepsilon_{\text{eq}}}{\varepsilon_{1/2}} - 1 \right) \right] + 0.5 + g_{\text{res}} \right),
\end{equation}
which lowers $G_c$ with accumulated plastic strain. Parameters: $G_{c0}$ initial toughness, $\ell_0$ regularization length, $g_{\text{res}}$ residual toughness ratio, and $\varepsilon_{1/2}$ the equivalent plastic strain at half-toughness. The equivalent plastic strain $\varepsilon_{\text{eq}}$ is derived from the PFC density gradients or an explicit plastic field, matching the code path in \texttt{PFCCoupling}.
Crack driving compares tensile elastic energy density with the critical fracture energy density $G_c/\ell_0$, so the numerical driving term uses $H \ell_0 /(G_c + \epsilon)$; this reintroduces the characteristic length to maintain a dimensionless ratio and avoid blow-up.

\paragraph{Plastic surrogates and anisotropic driving} Mechanical and PFC contributions are blended, with mechanical part driven by FCC slip RSS (orthogonal \{111\}$\langle 110\rangle$, sign-preserving, with yield threshold):
\begin{align}
\hat{\tau}_{\text{RSS}} &= \max_{m,n\in \{111\}\langle 110\rangle} |m\cdot(\varepsilon n)|, \qquad
\hat{g} = \frac{\sqrt{\sum_i (\partial_i \psi)^2}}{\max|\cdot|}, \\
\varepsilon_{\text{eq}} &= w_{\text{mech}}\,\frac{\hat{\tau}_{\text{RSS}}}{\max \hat{\tau}_{\text{RSS}}} + (1-w_{\text{mech}})\,\hat{g},\\
\varepsilon^p_{\text{flow}} &= \mathrm{flow}(\hat{\tau}_{\text{RSS}}-\tau_y)\,\mathrm{sign}(m\cdot\varepsilon n)\,\mathrm{sym}(m\otimes n),\\
\varepsilon_{\text{eq}}^{\text{acc}} &\leftarrow \varepsilon_{\text{eq}}^{\text{acc}} + \Delta t\,\gamma\,\varepsilon_{\text{eq}}, \qquad
\varepsilon^p \leftarrow \varepsilon^p + \Delta t\,\gamma\,\varepsilon^p_{\text{flow}},\\
\mathrm{flow}(x) &= \max\!\left(0, \frac{x}{1-\tau_y}\right)\, s_{\text{flow}},\\
\varepsilon_{\text{dir},i} &= w_{\text{mech}}\,[\varepsilon_{ii}]_+ + (1-w_{\text{mech}})\,\frac{|\partial_i \psi|}{\max_j |\partial_j \psi|},
\end{align}
where $w_{\text{mech}}\in[0,1]$ (default 0.9) and $[\cdot]_+$ clamps to tension. The loading-axis component $\varepsilon_{\text{dir},L}$ boosts both the history and crack driving:
\begin{equation}
H \leftarrow \max\!\big(H,\; f_{\text{elas}}^{+}(1+k_{\text{dir}}\varepsilon_{\text{dir},L})\big), \qquad
\mathcal{D} = \frac{H\,\ell_0}{G_c+\epsilon}(1-\phi)(1+k_{\text{dir}}\varepsilon_{\text{dir},L}),
\end{equation}
with $k_{\text{dir}}$ tuning anisotropy. When stress coupling is enabled, the chemical potential gets an additive term
\begin{equation}
\mu_{\text{extra}} = \alpha \frac{\sigma_{\text{vm}}}{\max |\sigma_{\text{vm}}|},
\end{equation}
so high von Mises regions accelerate PFC evolution. In mechanics, the accumulated plastic tensor $\varepsilon^p$ enters as an eigenstrain, so the stress is $C:(\varepsilon-\varepsilon^p)$; the accumulated scalar $\varepsilon_{\text{eq}}^{\text{acc}}$ degrades toughness $G_c(\varepsilon_{\text{eq}}^{\text{acc}})$ (optionally weakened on grain boundaries).

The PFC energy follows
\begin{equation}
f_{\text{PFC}} = \frac{\psi}{2} \left[r + (q_0^2 + \nabla^2)^2 \right]\psi + \frac{u}{4} \psi^4,
\end{equation}
with constraints (density conservation, grain fractions, plastic incompressibility) enforced through Lagrange multipliers in the modified TDGL formulation. Here $r$ is the undercooling parameter that sets the ordering tendency, $q_0$ is the principal reciprocal lattice magnitude, and $u$ penalizes large density excursions; the noise amplitude only seeds initial defects and does not enter the energy.

\section{Mechanical Equilibrium}
Small-strain equilibrium is imposed:
\begin{equation}
\nabla\cdot \boldsymbol{\sigma} = \mathbf{0}, \qquad \boldsymbol{\sigma} = C'(\mathbf{x}) : (\varepsilon - \varepsilon^p),
\end{equation}
with periodic in-plane boundaries and prescribed macroscopic strain $\varepsilon^{\text{macro}}$ in the loading direction. A conjugate-gradient linear operator solves for displacement $\mathbf{u}$, yielding total strain $\varepsilon = \nabla_s \mathbf{u} + \varepsilon^{\text{macro}}$. The rotated tensor $C'(\mathbf{x})$ comes from local orientation $R(\mathbf{x})$; $\varepsilon^p$ is either driven by PFC gradients or provided explicitly.

\section{Phase-Field Crystal Dynamics}
PFC density obeys conserved dynamics:
\begin{equation}
\frac{\partial \psi}{\partial t} = -\nabla^2 \mu, \qquad \mu = r\psi + (q_0^2 + \nabla^2)^2 \psi + u \psi^3,
\end{equation}
implemented spectrally with FFTs to evaluate $(q_0^2 - k^2)^2$ and advance $\psi$ semi-implicitly.

\section{Crack Evolution}
A history variable $H(\mathbf{x}) = \max_{t'\leq t} f_{\text{elas}}^{+}(\mathbf{x},t')$ accumulates tensile energy. Crack driving force combines energy and toughness:
\begin{equation}
\mathcal{D} = \frac{H\,\ell_0}{G_c + \epsilon}(1 - \phi),
\end{equation}
which compares the current tensile energy density to the critical fracture energy density $G_c / \ell_0$ (here $\ell_0 \approx 3\times10^{-6}$~m). The evolution law $\dot{\phi} = L_\phi \mathcal{D}$ is clamped to $[0,1]$, so cracks grow faster where energy is high and toughness is low.
Here $H$ stores the maximum tensile elastic energy density seen so far, $\epsilon$ is a small numerical residual stiffness, and $L_\phi$ is the kinetic coefficient for the crack field.

\section{Numerical Workflow}
\begin{enumerate}
    \item \textbf{Structure generation}: `Cu111StructureBuilder` constructs grids (32$\times$32$\times$16, spacing $5\times10^{-7}$~m) with $[111]$ orientation and random defects for testing; EBSD data can replace the synthetic mask in future runs. Defect seeding uses weighted/random seeds (and optional line defects) projected via anisotropic Gaussians to $\psi$/crack/plastic; optional notch boxes pre-set high crack values as starters.
    \item \textbf{Alternating solver}: The cyclic load path is triangular $0\rightarrow +\varepsilon_{\max}\rightarrow 0 \rightarrow -\varepsilon_{\max}\rightarrow 0$, each segment subdivided into $N_{\text{seg}}$ steps. At every increment: solve mechanical equilibrium (CG); blend mechanical/PFC plastic surrogates and relax plastic scalar + directional vector; update crack with anisotropic boost; advance PFC with optional stress-coupled $\mu_{\text{extra}}$.
    \item \textbf{Macro strain with Poisson contraction}: Applied strain is $(\varepsilon, -\nu \varepsilon, -\nu \varepsilon)$ (default $\nu=0.34$ for Cu); the same macro strain is used in PFC $k$-update, VTK warp, and LAMMPS dump.
    \item \textbf{Toughness scaling}: A scalar $s_{\mathrm{Gc}}$ rescales $g_c$ and $g_{\mathrm{res}}$ (e.g., $s<1$ to promote cracking), leaving $\ell_0,k,\epsilon_{1/2}$ unchanged.
    \item \textbf{Outputs}: Every few steps dump VTK frames (deformed coordinates) and after each cycle dump LAMMPS and optional `.data` snapshots. Scalar/vector fields include crack, plastic (scalar + directional), $\psi$, stress (tensor and von Mises), displacement, and masks (defect/line/notch). CSV logs store per-cycle energy, mean crack/plastic, crack increment, plastic range, and (optionally) a stepwise macro stress--strain curve $\{\varepsilon, \bar{\sigma}_{xx}, \bar{\sigma}_{yy}, \bar{\sigma}_{zz}, \bar{\sigma}_{\mathrm{vm}}\}$; runs stop early if $\langle\phi\rangle$ exceeds a failure threshold.
    \item \textbf{Fatigue metrics}: Fit crack increment per cycle vs cycle count (Paris-like slope) and plastic strain range vs $2N$ (Coffin--Manson slope) from the logged series.
\end{enumerate}

\section{Defect Seeding and Initialization (New)}
\paragraph{Weighted random and deterministic line seeds} A defect configuration specifies: density $\rho_d$, region bounds $\Omega_d$, type probabilities $p_k$ (vacancy/interstitial/dislocation), optional weight field $w(\mathbf{x})$, and optional line segments $\{\Gamma_\ell\}$. The total target seeds is $N=\min(\rho_d |\Omega_d|, N_{\max})$. We first discretize each line $\Gamma_\ell$ into points $\mathbf{x}_i$ with spacing $\Delta s$ and mark them as dislocation seeds (anisotropic kernel along the line); the remaining $N_{\text{rem}}$ seeds are drawn from a categorical distribution over types and sampled positions:
\begin{equation}
P(\mathbf{x}) = \frac{w(\mathbf{x})\,\chi_{\Omega_d}(\mathbf{x})}{\int_{\Omega_d} w(\mathbf{x})\,d\mathbf{x}}, \quad w(\mathbf{x}) = \begin{cases}0,&w< w_{\text{th}}\\ w,&\text{else}\end{cases},
\end{equation}
falling back to uniform if $w$ is absent. Directions are random or sampled from FCC slip systems $\langle 110\rangle$.

\paragraph{Gaussian projection to fields} Each seed contributes a 3D anisotropic Gaussian kernel to $\psi$, crack, and plastic fields. For a seed at $\mathbf{x}_0$ with kernel widths $\boldsymbol{\sigma}$ and weight $A$:
\begin{equation}
g(\mathbf{x}) = A \exp\!\left[-\sum_{i=1}^3 \frac{(x_i-x_{0,i})^2}{2\sigma_i^2}\right], \quad
\psi \leftarrow \psi + g,\ \phi \leftarrow \phi + A_\phi g,\ \varepsilon^p \leftarrow \varepsilon^p + A_p g.
\end{equation}
Vacancy uses negative $A$, interstitial/dislocation use positive; line seeds use elongated $\boldsymbol{\sigma}$ along the line direction. After accumulation, $\psi$ is rescaled to a target mean/peak ($\bar{\psi}\rightarrow \bar{\psi}_{\text{tgt}}$, $\max|\psi|\rightarrow \psi_{\text{peak}}$), and fields are clipped to $[0,1]$. Masks \texttt{defect\_mask} (all seeds) and \texttt{line\_mask} (line seeds) are written for visualization.

\paragraph{EBSD orientation ingestion} The builder accepts an external orientation map $R(\mathbf{x})$ (e.g., EBSD) to replace the synthetic $[111]$ field, ensuring stiffness rotation and subsequent mechanics follow the measured texture.

\paragraph{Per-step VTK sequence} A smoke test outputs VTK per time step with deformed coordinates (\texttt{displacement\_total}) to allow ParaView time-series playback; final snapshots include both deformed and undeformed grids for easy comparison.

\section{Outlook}
The framework replicates Fig.~1--4 behavior from Xu \emph{et al.}\cite{paper} at the single-crystal scale. Next steps include (i) ingesting actual EBSD orientation/defect maps, (ii) swapping the CG equilibrium for FFT-based Lippmann--Schwinger or FEM for faster convergence and non-periodic BCs, and (iii) embedding full crystal-plasticity slip kinetics and validated fracture history laws.

\begin{thebibliography}{1}
\bibitem{paper} F. Xue \emph{et al.}, ``Phase-field framework with constraints and its applications to ductile fracture in polycrystals and fatigue,'' \emph{npj Computational Materials} 8, 18 (2022).
\end{thebibliography}

\section{Recent Code Updates (Summary)}
\begin{itemize}
    \item 塑性与方向场：新增 \texttt{plastic\_measures} 计算 $\varepsilon_{\text{eq}}$ 与方向分量，机械 von Mises/轴向应变与 PFC 梯度按权重混合，默认机械占比 0.9。
    \item 方向性裂纹驱动：加载轴塑性正分量放大历史能量和驱动力（\texttt{dir\_coupling}=0.8），抑制非加载方向裂纹。
    \item 应力耦合 PFC：von Mises 应力归一化后作为 $\mu_\text{extra}$ 加入化学势，提升高应力区的晶格演化响应。
    \item 求解与载荷：采用对称三角循环载荷（0$\rightarrow$+$\varepsilon_{\max}\rightarrow$0$\rightarrow$-$\varepsilon_{\max}\rightarrow$0），可设置失败阈值提前终止；跟踪每周期塑性范围 $\Delta\varepsilon_p$ 和裂纹增量。
    \item VTK/LAMMPS 输出扩展：输出塑性方向分量、应力张量/ von Mises，附带归一化场（crack\_norm, plastic\_vec\_norm, stress\_vm\_norm, disp\_total\_norm 等）便于统一色标；VTK 现默认二进制 STRUCTURED\_GRID（变形坐标），避免 ASCII 读入报错，LAMMPS 去重了多余的 \texttt{plastic\_vec} 列。
    \item 结构缺陷默认幅值调至 0.12，PFC 截断默认 $|\psi|\leq 1.2$，方便在高载下保持稳定。
\end{itemize}

\section{Practical Simulation Recipe}
\begin{enumerate}
    \item \textbf{结构初始化}：用 \texttt{Cu111StructureBuilder} 构造网格（默认 [111] 单晶，或提供 grain labels/orientation map）。缺陷幅值默认 0.12，噪声注入 $\psi$；可直接替换为 EBSD 导出的取向场。
    \item \textbf{求解循环}：\texttt{run\_virtual\_cycles} 采用三角载荷 0$\rightarrow$+$\varepsilon_\text{max}\rightarrow$0$\rightarrow$-$\varepsilon_\text{max}\rightarrow$0，默认每段 50 步，可多周期。每步执行：力学求解（CG）；塑性标量+方向分量松弛；裂纹更新（方向增益）；PFC 更新（含应力耦合 $\mu_{\text{extra}}$）。
    \item \textbf{参数提示}：\texttt{mech\_plastic\_weight}=0.9，\texttt{dir\_coupling}=0.8，\texttt{plastic\_relax}=0.2，\texttt{clip}=1.2，\texttt{failure\_threshold}=0.98；必要时调节应力耦合系数（\texttt{mu\_extra\_from\_stress}）或加载轴 \texttt{load\_axis}。
    \item \textbf{输出与监控}：每 5 步 VTK 帧（\texttt{anim\_frame\_*.vtk}）和每周期 LAMMPS dump（\texttt{virtual\_cycle\_*.lammpstrj}）；CSV 日志记录能量、裂纹均值/增量、塑性均值/范围，裂纹均值超阈值将提前停止。
\end{enumerate}

\section{Field/Output Guide}
\begin{itemize}
    \item \textbf{VTK frame files} (\texttt{sim/tests/virtual\_cycle\_vtk/anim\_frame\_*.vtk}): 二进制 STRUCTURED\_GRID（若启用变形坐标），含下列标量/向量；加载时关闭 per-timestep rescale，固定色标（裂纹 0--0.3 或 0--1，塑性/应力 0--1）。
    \item $\psi$: PFC 密度偏差，红/蓝表示高/低密度或缺陷核；受应力耦合影响会在高应力区演化。
    \item \texttt{crack}, \texttt{crack\_norm}, \texttt{crack\_clamp03}: 裂纹变量（0--1），clamp03 便于观察萌生阶段，norm 便于固定色标。
    \item \texttt{plastic} (标量), \texttt{plastic\_vec} (方向): 来自机械应变 + PFC 梯度；\texttt{plastic\_vec\_x\_norm} 可直接用于加载方向色标，\texttt{plastic\_vec\_mag} 为模值。
    \item \texttt{stress\_vm}, \texttt{stress\_vm\_norm}: von Mises 应力，可用于判断高应力区；亦用于 $\mu_\text{extra}$。
    \item \texttt{displacement}, \texttt{displacement\_total}: 微观/总位移；\texttt{disp\_total\_norm} 提供归一化幅值，避免色标跳变。
    \item \textbf{LAMMPS dump} (\texttt{sim/tests/virtual\_cycle\_0001.lammpstrj}): 列顺序为 \texttt{id, type, x, y, z, crack, plastic, psi, plastic\_xx, plastic\_yy, plastic\_zz, stress\_vm, stress\_xx, stress\_yy, stress\_zz}（若场缺失则省略），坐标包含宏+微观位移。
    \item \textbf{CSV 日志} (\texttt{sim/tests/virtual\_cycle.csv}): \texttt{cycle}（循环编号）、\texttt{energy}（总能量）、\texttt{crack\_mean}（裂纹均值）、\texttt{plastic\_mean}（塑性均值）、\texttt{crack\_delta}（该周期裂纹增量）、\texttt{plastic\_range}（该周期塑性范围 $\Delta\varepsilon_p$）。
    \item 归一化场均在 VTK 中提供，便于在 ParaView 固定色标（range 0--1 或 0--0.3）播放时间序列。
\end{itemize}

\section{实验流程与数据对应}
\begin{itemize}
    \item \textbf{载荷与网格}: 默认单晶网格 $32\times32\times16$，缺陷幅值 0.12；三角循环载荷（0$\rightarrow$+$\varepsilon_{\max}\rightarrow$0$\rightarrow$-$\varepsilon_{\max}\rightarrow$0），每段 50 步，\texttt{max\_strain}=0.08，可设置 \texttt{failure\_threshold}（0.98）提前终止。
    \item \textbf{周期统计}: \texttt{energy} 为周期末总能量；\texttt{crack\_mean} 代表损伤程度（接近 1 表示贯穿）；\texttt{crack\_delta} 是相邻周期裂纹增量，用于 Paris 拟合；\texttt{plastic\_mean} 为塑性均值；\texttt{plastic\_range} 为该周期塑性最大-最小（$\Delta\varepsilon_p$ 替代量，用于 Coffin--Manson 拟合）。
    \item \textbf{可视化要点}: 在 ParaView/Ovito 打开 VTK 或 LAMMPS dump 时，优先查看 \texttt{crack\_clamp03}（萌生）、\texttt{plastic\_vec\_x\_norm}（加载方向塑性）、\texttt{stress\_vm\_norm}（等效应力）、\texttt{disp\_total\_norm}（总位移）。固定色标 0--1 便于跨帧比较。
    \item \textbf{Coffin--Manson/Paris 指标}: 拟合对数斜率 $\mathrm{d}\log(\Delta\varepsilon_p)/\mathrm{d}\log(2N)$ 得到 Coffin--Manson；拟合 $\log(\Delta\phi)/\log(N)$ 得到 Paris-like 系数，均由 \texttt{sim/tests/virtual\_cycle.py} 自动计算并打印。
\end{itemize}

\section{当前结果快照（截至 2026-02-09）}
\begin{itemize}
    \item \textbf{D1 全量门禁通过}: \texttt{\detokenize{sim/tests/regress_runs/2026-02-09/d1_full_gate/summary.json}}，总状态 \texttt{passed=true}。验收矩阵为 \texttt{phase2\_full=true}、\texttt{multi\_align\_full=5/5}、\texttt{seed\_robustness=6/6}。
    \item \textbf{D2 裂纹局部化 + 能量密度门禁通过}: \texttt{\detokenize{sim/tests/regress_runs/2026-02-09/d2_localization_energy/summary.json}}，\texttt{cycles\_completed=6}，\texttt{crack\_delta\_total=2.396875}，\texttt{crack\_localization\_index\_peak=5.1072}，并输出 6 个能量相关 VTK 字段。
    \item \textbf{D3 多物理矩阵通过}: \texttt{\detokenize{sim/tests/regress_runs/2026-02-09/d3_multiphysics_matrix_smoke/summary.json}}，\texttt{case\_total=3}，\texttt{passed\_count=3}；包含 2 个正向裂纹触发 case 与 1 个负对照抑制 case。
    \item \textbf{Release D2+D3 quick+seedfull 通过}: \texttt{\detokenize{sim/tests/regress_runs/2026-02-09/release_d2_d3_quick_seedfull/bundle_summary.json}}，\texttt{acceptance.d2\_localization=true}、\texttt{acceptance.d3\_matrix=true}，并且两批 seed（41--46）均 \texttt{all\_seed\_gate\_passed=true}。
\end{itemize}

\section{全量基准测试矩阵（脚本级）}
\begin{itemize}
    \item \texttt{\detokenize{sim/tests/regress_microstrain.py}}：线弹性微应变比值回归；最新通过产物 \texttt{\detokenize{sim/tests/regress_runs/2026-02-05/microstrain/summary.json}}（\texttt{stress\_ratio\,$\approx$\,2.0}）。
    \item \texttt{\detokenize{sim/tests/regress_gnd.py}}：GND/Nye 线性响应回归；最新通过产物 \texttt{\detokenize{sim/tests/regress_runs/2026-02-05/gnd/summary.json}}（双梯度比值约 2）。
    \item \texttt{\detokenize{sim/tests/regress_gnd_cycle.py}}：低幅循环 GND 增长回归；最新通过产物 \texttt{\detokenize{sim/tests/regress_runs/2026-02-05/gnd_cycle/summary.json}}（\texttt{gnd\_growth>0}、\texttt{accum\_growth>0}）。
    \item \texttt{\detokenize{sim/tests/regress_bc_crack.py}}：边界/裂纹小规模回归；最新通过产物 \texttt{\detokenize{sim/tests/regress_runs/2026-02-05/bc_crack/summary.json}}。
    \item \texttt{\detokenize{sim/tests/regress_bc_crack_large.py}}：大网格边界/裂纹回归；最新通过产物 \texttt{\detokenize{sim/tests/regress_runs/2026-02-05/bc_crack_large/summary.json}}。
    \item \texttt{\detokenize{sim/tests/regress_bc_crack_micron.py}}：微米尺度缩放回归；最新通过产物 \texttt{\detokenize{sim/tests/regress_runs/2026-02-05/bc_crack_micron/summary.json}}。
    \item \texttt{\detokenize{sim/tests/regress_all.py}}：边界裂纹三件套统一入口；最新通过产物 \texttt{\detokenize{sim/tests/regress_runs/2026-02-06/phase1_baseline/boundary_crack/summary.json}}（small/large/micron 全通过）。
    \item \texttt{\detokenize{sim/tests/run_phase1_suite.py}}：Phase-1 一键回归；最新通过产物 \texttt{\detokenize{sim/tests/regress_runs/2026-02-06/phase1_suite_venvcheck/summary.json}}（fresh venv 复现通过）。
    \item \texttt{\detokenize{sim/tests/regress_phase2.py}}：Phase-2 综合门禁编排；最新通过产物 \texttt{\detokenize{sim/tests/regress_runs/2026-02-09/release_d2_d3_quick_seedfull/phase2_gate/summary.json}}（exp + D2 同时开启）。
    \item \texttt{\detokenize{sim/tests/regress_exp_alignment.py}}：单工况实验对齐回归；最新通过产物 \texttt{\detokenize{sim/tests/regress_runs/2026-02-09/release_d2_d3_quick_seedfull/phase2_gate/exp_alignment/summary.json}}（\texttt{rmse\_tau=28.45 MPa}，\texttt{rmse\_gamma=0.003889}）。
    \item \texttt{\detokenize{sim/tests/regress_energy_consistency.py}}：能量/裂纹/塑性一致性门禁；最新通过产物 \texttt{\detokenize{sim/tests/regress_runs/2026-02-08/energy_gate_smoke/summary.json}}（无能量下降和反转计数）。
    \item \texttt{\detokenize{sim/tests/regress_exp_alignment_multi.py}}：多工况实验对齐回归；最新通过产物 \texttt{\detokenize{sim/tests/regress_runs/2026-02-08/exp_alignment_multi_week8_smoke/summary.json}}（5/5 条件通过）。
    \item \texttt{\detokenize{sim/tests/repeat_crack_onset_seeds.py}}：seed 稳健性批次回归；最新通过产物 \texttt{\detokenize{sim/tests/regress_runs/2026-02-09/d1_full_gate/seed_batch_1/summary.json}} 与 \texttt{\detokenize{sim/tests/regress_runs/2026-02-09/d1_full_gate/seed_batch_2/summary.json}}（两批均 3/3）。
    \item \texttt{\detokenize{sim/tests/regress_d2_localization_energy.py}}：D2 裂纹局部化触发与能量场输出回归；最新通过产物 \texttt{\detokenize{sim/tests/regress_runs/2026-02-09/d2_localization_energy/summary.json}}。
    \item \texttt{\detokenize{sim/tests/regress_d3_multiphysics_matrix.py}}：D3 正/负多物理矩阵回归；最新通过产物 \texttt{\detokenize{sim/tests/regress_runs/2026-02-09/d3_multiphysics_matrix_smoke/summary.json}}（3/3 通过）。
    \item \texttt{\detokenize{sim/tests/run_d1_full_gate.py}}：D1 全量编排入口；最新通过产物 \texttt{\detokenize{sim/tests/regress_runs/2026-02-09/d1_full_gate/summary.json}}（验收矩阵全绿）。
    \item \texttt{\detokenize{sim/tests/run_release_baseline_week4.py}}：release 编排（Phase-2 + D2 + D3 + seed）；最新通过产物 \texttt{\detokenize{sim/tests/regress_runs/2026-02-09/release_d2_d3_quick_seedfull/bundle_summary.json}}。
    \item \texttt{\detokenize{sim/tests/run_ci_smoke.py}}：CI 最小门禁链路；最新通过产物 \texttt{\detokenize{sim/tests/regress_runs/2026-02-08/ci_smoke_local_verify_fix/summary.json}}（3 个子任务全部通过）。
    \item \texttt{\detokenize{sim/tests/run_seed_robustness_20.py}}：20-seed 分批运行模板；最新通过产物 \texttt{\detokenize{sim/tests/regress_runs/2026-02-08/seed_robustness_20_smoke_pair/bundle_summary.json}}（smoke 模式通过）。
    \item \texttt{\detokenize{sim/tests/summarize_seed_robustness_ci.py}}：seed CI 统计聚合；最新通过产物 \texttt{\detokenize{sim/tests/regress_runs/2026-02-08/seed_ci_summary_smoke/summary.json}}（\texttt{all\_seed\_gate\_passed=true}）。
\end{itemize}

\section{尚需完善与下一步}
\begin{itemize}
    \item \textbf{D2/D3 全量化}: 当前已有 quick+seedfull 基线，下一步应固定 full profile 常态门禁，并保证与 quick 指标一致。
    \item \textbf{Seed 统计扩容}: 将 seed 稳健性从 6 seeds 扩展到 20 seeds full-case，产出稳定的 Wilson 区间与聚合报告。
    \item \textbf{D3 负对照增强}: 增加更多“无缺口/弱驱动/不同晶向”负样本，降低误触发风险并提升外推稳健性。
    \item \textbf{长周期验证}: 建立 \(>10\) cycles 的 D2+D3 联合回归，补齐能量密度与裂纹局部化随周期演化证据链。
    \item \textbf{发布摘要自动化}: 在现有 \texttt{bundle\_summary.json} 基础上生成统一 markdown/tex 摘要，减少手工整理成本。
\end{itemize}

\end{document}
